\subsection{Analisi Within Country e Test dei Segni}\label{sec:methods_within}
%Per analizzare l'home advantage è stato effettuato un confronto within-country:
%la performance di ogni Paese che ha ospitato almeno una volta è stata
%confrontata con la media delle performance dello stesso Paese quando non ha
%ospitato. 
%Obiettivo di tale analisi è verificare se esiste un vantaggio sistematico quando un Paese ospita le Olimpiadi.

%Calcolati i delta, ovvero le differenze fra la percentuale di medaglie vinte quando un Paese è ospitante e
%la media delle percentuali di medaglie vinte quando non lo è, è stato effettuato un Test dei Segni su tali delta.
%Il Test dei Segni è un test statistico non parametrico utilizzato per
%determinare se esiste una differenza sistematica tra coppie di osservazioni
%correlate.
%Tale test è statisticamente meno potente rispetto ad alternative come il t Test
%o il test di Wilcoxon.
%Tuttavia, questi due test richiedono che i dati provengano da distribuzioni normali o simmetriche,
%requisiti che abbiamo verificato che i nostri dati non rispettano.
